\chapter*{全书符号、单位与约定}
\addcontentsline{toc}{chapter}{全书符号、单位与约定}

除特别说明外，全书采用 SI 制并显式保留 $\hbar,c,\epsilon_0,\mu_0$；微分、虚数单位和指数分别统一写作 $\dd,\ii,\ee$。空间矢量使用统一矢量字体，矩阵与二阶以上线性算子按各章定义使用固定字体。Fourier 变换、连续态归一化与 Green 函数边界条件在首次使用时明确给出，后续不得在同一模块内改变号数或归一化。

本书固定以下跨章注册：统计物理中的粒子数密度记为 $\varrho$；薄板三维质量密度记为 $\rho_{\rm m}$，面质量密度记为 $\mu_A$；电磁学中的电荷密度记为 $\rho_{\rm e}$；开放量子系统的密度算符保留 $\rho$。辐射与高维散射中的 $\dd\Omega$ 专指立体角元，开放系统频谱积分若需独立哑频率则使用 $\varpi$，而 $\omega$ 保留为物理角频率或 Bohr 频率。高维中心势中的 $\mu$ 表示约化质量，不与薄板的 $\mu_A$ 混用。

Ising 部分把各向异性对角相关的低温参数固定为 $\kappa=[\sinh(2K_h)\sinh(2K_v)]^{-1}$；簇 expansion 的无量纲系数固定为 $b_n$，三体 connected canonical cumulant 记为 $C_3$。同一物理对象在全书只保留一个主要符号；同一章节内不允许一个符号同时承担不同数学类型。传统上不可避免的局部字母复用，只能发生在明确分隔的模块中，并必须在首次出现处重新定义类型、单位和物理意义。
